% Ad Quintum Fratrem 1.2 (ad-quintum-fratrem-01-02)
% Source: https://raw.githubusercontent.com/PerseusDL/canonical-latinLit/master/data/phi0474/phi058/phi0474.phi058.perseus-lat1.xml
% Fetched by scripts/fetch_latin.py

% Dateline (Perseus): Scr. Romae inter va vi K. Nov. et iv Id. Dec. a. 695 (59).

\ciceroLetterOpener{MARCVS QVINTO FRATRI}

\ciceroSection{1} Statius ad me venit a. d. viii K. Novembr. eius adventus, quod ita scripsisti direptum iri te a tuis dum is abesset, molestus mihi fuit ; quod autem exspectationem tui concursumque eum qui erat futurus si una tecum decederet neque antea visus esset sustulit, id mihi non incommode visum est accidisse. exhaustus est enim sermo hominum et multae emissae iam eius modi voces: ' ἀλλὰ αἰεί τινα φῶτα μέγαν ,' quae te absente confecta esse laetor. quod autem idcirco a te missus est mihi ut se purgaret, id necesse minime fuit.

\ciceroSection{2} primum enim numquam ille mihi fuit suspectus neque ego, quae ad te de illo scripsi, scripsi meo iudicio, sed cum ratio salusque omnium nostrum qui ad rem publicam accedimus non veritate solum sed etiam fama niteretur, sermones ad te aliorum semper, non mea iudicia perscripsi. qui quidem quam frequentes essent et quam graves adventu suo Statius ipse cognovit ; etenim intervenit non nullorum querelis quae apud me de illo ipso, habebantur, et sentire potuit sermones iniquorum in suum potissimum nomen erumpere.

\ciceroSection{3} quod autem me maxime movere solebat, cum audiebam illum plus apud te posse quam gravitas istius aetatis, imperi, prudentiae postularet quam multos enim mecum egisse putas ut se Statio commendarem, quam multa autem ipsum ἀφελῶσ mecum in sermone ita posuisse : ' id mihi non placuit ; monui, suasi, deterrui'! quibus in rebus etiam si fidelitas summa est, (quod prorsus credo quoniam tu ita iudicas), tamen species ipsa tam gratiosi liberti aut servi dignitatem habere nullam potest. atque hoc sic habeto (nihil enim nec temere dicere nec astute reticere debeo), materiam omnem sermonum eorum qui de te detrahere vellent Statium dedisse; antea tantum intellegi potuisse, iratos tuae severitati esse non nullos, hoc manumisso iratis quod loquerentur non defuisse.

\ciceroSection{4} nunc respondebo ad eas epistulas quas mihi reddidit L. Caesius, quoi quoniam ita te velle intellego nullo loco deero ; quarum altera est de Blaundeno Zeuxide, quem scribis certissimum matricidam tibi a me intime commendari. qua de re et de hoc genere toto, ne forte me in Graecos tam ambitiosum factum esse mirere pauca cognosce. ego cum Graecorum querelas nimium valere sentirem propter hominum ingenia ad fallendum parata, quoscumque de te queri audivi quacumque potui ratione placavi. primum Dionysopolitas qui erant inimicissimi mei lenivi ; is quorum principem Hermippum non solum sermone meo sed etiam familiaritate devinxi. ego Apamensem Hephaestium, ego levissimum hominem, Megaristum Antandrium, ego Niciam Smyrnaeum, ego nugas maximas omni mea comitate complexus sum, Nymphontem etiam Colophonium. quae feci omnia, non quo me aut hi homines aut tota natio delectaret ; pertaesum est levitatis, adsentationis, animorum non officiis sed temporibus servientium.

\ciceroSection{5} sed ut ad Zeuxim revertar, cum is de M. Cascelli sermone secum habito quae tu scribis ea ipsa loqueretur, obstiti eius sermoni et hominem in familiaritatem recept tua autem quae fuerit cupiditas tanta nescio, quod scribis cupisse te, quoniam Smyrnae duos Mysos insuisses in culleum, simile in superiore parte provinciae edere exemplum severitatis tuae et idcirco Zeuxim elicere omni ratione voluisse, ultro quem adductum in iudicium fortasse dimitti non oportuerat, conquiri vero et elici blanditus, ut tu scribis, ad iudicium necesse non fuit, eum praesertim hominem quem ego et ex suis civibus et ex multis aliis cotidie magis cognosco nobiliorem esse prope quam civitatem suam.

\ciceroSection{6} at enim Graecis solis indulgeo. quid ? L. Caecilium nonne omni ratione placavi? quem hominem, qua ira, quo spiritu! quem denique praeter Tuscenium cuius causa sanari non potest non mitigavi? ecce supra caput homo levis ac sordidus sed tamen equestri censu, Catienus. etiam is lenietur. cuius tu in patrem quod fuisti asperior non reprehendo ; certo enim scio te fecisse cum causa. sed quid opus fuit eius modi litteris quas ad ipsum misisti, 'illum crucem sibi ipsum constituere, ex qua tu eum ante detraxisses ; te curaturum fumo ut combureretur plaudente tota provincia'? quid vero ad C. Fabium nescio quem (nam eam quoque epistulam T. Catienus circumgestat), 'renuntiari tibi Licinium plagiarium cum suo pullo milvino tributa exigere'? deinde rogas Fabium ut et patrem et filium vivos comburat si possit ; si minus, ad te mittat uti iudicio comburantur. eae litterae abs te per iocum missae ad C. Fabium, si modo sunt tuae, cum leguntur, invidiosam atrocitatem verborum habent.

\ciceroSection{7} ac si omnium mearum praecepta litterarum repetes, intelleges esse nihil a me nisi orationis acerbitatem et iracundiam et, si forte, raro litterarum missarum indiligentiam reprehensam. quibus quidem in rebus si apud te plus auctoritas mea quam tua sive natura paulo acrior sive quaedam dulcedo iracundiae sive dicendi sal facetiaeque valuissent, nihil sane esset quod nos paeniteret. et mediocri me dolore putas adfici cum audiam qua sit existimatione C. Vergilius, qua tuus vicinus, C. Octavius? nam si te interioribus vicinis tuis, Ciliciensi et Syriaco, anteponis, valde magni facis! atque is dolor est quod, cum ii quos nominavi te innocentia non vincant, vincunt tamen artificio benevolentiae conligendae, qui neque Cyrum Xenophontis neque Agesilaum noverint, quorum regum summo imperio nemo umquam verbum ullum asperius audivit.

\ciceroSection{8} sed haec a principio tibi praecipiens quantum profecerim non ignoro ; nunc tamen decedens, id quod mihi iam facere videris, relinque, quaeso, quam iucundissimam memoriam tui. successorem habes perblandum ; cetera valde illius adventu tua requirentur. in litteris mittendis, ut saepe ad te scripsi, nimium te exorabilem praebuisti. tolle omnis, si potes, iniquas, tolle inusitatas, tolle contrarias. Statius mihi narravit scriptas ad te solere adferri, ab se legi, et si iniquae sint fieri te certiorem ; ante quam vero ipse ad te venisset, nullum delectum litterarum fuisse ; ex eo esse volumina selectarum epistularum quae reprehendi solerent.

\ciceroSection{9} hoc de genere nihil te nunc quidem moneo (sero est enim) ac scire potes multa me varie diligenterque monuisse ; illud tamen quod Theopompo mandavi cum essem admonitus ab ipso, vide per homines amantis tui, quod est facile, ut haec genera tollantur epistularum primum iniquarum, deinde contrariarum, tum absurde et inusitate scriptarum, postremo in aliquem contumeliosarum. atque ego haec tam esse quam audio non puto, et si sunt occupationibus tuis minus animadversa nunc perspice et purga. legi epistulam, quam ipse scripsisse Sulla nomenclator dictus est, non probandam, legi non nullas iracundas.

\ciceroSection{10} sed tempore ipso de epistulis. nam cum hanc paginam tenerem L. Flavius, praetor designatus, ad me venit, homo mihi valde familiaris. is mihi te ad procuratores suos litteras misisse, quae mihi visae sunt iniquissimae, ne quid de bonis, quae L. Octavi Nasonis fuissent cui L. Flavius heres est, deminuerent, ante quam C. Fundanio pecuniam solvissent, itemque misisse ad Apollonidensis ne de bonis, quae Octavi fuissent, deminui paterentur, prius quam Fundanio debitum solutum esset. haec mihi veri similia non videnr0 tur ; sunt enim a prudentia tua remotissima. ne deminuat heres? quid , si infitiatur? quid, si omnino non debet? quid? praetor solet iudicare deberi? quid ? ego Fundanio non cupio, non amicus sum, non misericordia moveor? nemo magis ; sed via iuris eius modi est quibusdam in rebus ut nihil sit loci gratiae. atque ita mihi dicebat Flavius scriptum in ea epistula quam tuam esse dicebat, te aut quasi amicis tuis gratias acturum aut quasi inimicis incommoda laturum.

\ciceroSection{11} quid multa? ferebat graviter et vehementer mecum querebatur orabatque ut ad te quam diligentissime scriberem. quod facio et te prorsus vehementer etiam atque etiam rogo ut te procuratoribus Flavi remittas de deminuendo et Apollonidensibus ne quid praescribas, quod contra Flavium sit, amplius. et Flavi causa et scilicet Pompei facies omnia. nolo medius fidius ex tua iniuria in illum tibi liberalem me videri, sed et te oro ut tu ipse auctoritatem et monumentum aliquod decreti aut litterarum tuarum relinquas quod sit ad Flavi rem et ad causam accommodatum ; fert enim graviter homo et mei observantissimus et sui iuris dignitatisque retinens se apud te neque amicitia nec iure valuisse ; et, ut opinor, Flavi aliquando rem et Pompeius et Caesar tibi commendarunt et ipse ad te scripserat Flavius et ego certe. qua re, si ulla res est quam tibi me petente faciendam putes, haec ea sit. si me amas, cura, elabora, perfice ut Flavius et tibi et mihi quam maximas gratias agat. hoc te ita rogo ut maiore studio rogare non possim.

\ciceroSection{12} quod ad me de Hermia scribis mihi me hercule valde molestum fuit. Litteras ad te parum fraterne scripseram ; quas oratione Diodoti, Luculli liberti, commotus, de pactione statim quod audieram, iracundius scripseram et revocare cupiebam. huic tu epistulae non fraterne scriptae fraterne debes ignoscere.

\ciceroSection{13} de Censorino, Antonio, cassus , Scaevola te ab iis diligi , ut scribis, vehementer gaudeo. cetera fuerunt in eadem epistula graviora quam vellem, ' ὀρθὰν τὰν ναῦν ,' et ' ἅπαξ θανεῖν .' maiora ista erunt ; meae obiurgationes fuerunt amoris plenissimae † quae sunt non nulla, sed tamen mediocria et parva potius. ego te numquam ulla in re dignum minima reprehensione putassem, cum te sanctissime gereres, nisi inimicos multos haberemus. quae ad te aliqua cum monitione aut obiurgatione scripsi, scripsi propter diligentiam cautionis meae, in qua et maneo et manebo et idem ut facias non desistam rogare.

\ciceroSection{14} Attalus Hypaepenus mecum egit ut se ne impedires quo minus quod ad Q. †Publiceni† statuam decretum est erogaretur. quod ego te et rogo et admoneo ne talis viri tamque nostri necessari honorem minui per te aut impediri velis. praeterea Aesopi, nostri tragoedi familiaris, Licinus servus tibi notus aufugit. is Athenis apud Patronem Epicureum pro libero fuit ; inde in Asiam venit. postea Plato quidam Sardianus, Epicureus, qui Athenis solet esse multum et qui tum Athenis fuerat cum Licinius eo venisset, cum eum fugitivum esse postea ex Aesopi litteris cognosset, hominem comprehendit et in custodiam Ephesi tradidit, sed in publicam an in pistrinum non satis ex litteris eius intellegere potuimus. tu , quoquo modo est, quoniam Ephesi est, hominem investiges velim summaque diligentia vel tecum deducas. noli spectare quanti homo sit ; parvi enim preti est qui tam nihili sit. sed tanto dolore Aesopus est adfectus propter servi scelus et audaciam ut nihil ei gratius facere possis quam si illum per te reci perarit.

\ciceroSection{15} nunc ea cognosce quae maxime exoptas. rem publicam funditus amisimus, adeo ut Cato, adulescens nullius consili sed tamen civis Romanus et Cato, vix vivus effugerit quod, cum Gabinium de ambitu vellet postulare neque praetores diebus aliquot adiri possent vel potestatem sui facerent, in contionem escendit et Pompeium 'privatum dictatorem' appellavit. propius nihil est factum quam ut occideretur. ex hoc qui sit status totius rei publicae videre potes.

\ciceroSection{16} nostrae tamen causae non videntur homines defuturi ; mirandum in modum profitentur, offerunt se, pollicentur. equidem cum spe sum maxima tum maiore etiam animo, spe, superiores fore nos, animo, ut in hac re publica ne casum quidem ullum pertimescam sed tamen se res sic habet.: si diem nobis dixerit, tota Italia concurret, ut multiplicata gloria discedamus ; sin autem vi agere conabitur, spero fore studiis non solum amicorum sed etiam alienorum ut vi resistamus. omnes et se et suos amicos, clientis libertos, servos, pecunias denique suas pollicentur. nostra antiqua manus bonorum ardet studio nostri atque amore. si qui antea aut alieniores fuerant aut languidiores, nunc horum regum odio se cum bonis coniungunt. Pompeius omnia pollicetur et Caesar ; quibus ego ita credo ut nihil de mea comparatione deminuam. tribuni pl. designati sunt nobis amici ; consules se optime ostendunt ; praetores habemus amicissimos et acerrimos civis, Domitium, Nigidium, Memmium, Lentulum ; bonos etiam alios singularis. qua re magnum fac animum habeas et spem bonam de singulis tamen rebus quae cotidie gerantur faciam te crebro certiorem.
